%%%%%%%%%%%%%%%%%%%%%%%%%%%%%%%%%%%%%%%%%%%%%%%%%%%%%%%
% NOTAS TEORIA DE TRANSPORTE OPTIMO
%
% Contains ?? figures. 
% File started: ??
% First circulated version: ??
% Last version:
% Last modification: -
%%%%%%%%%%%%%%%%%%%%%%%%%%%%%%%%%%%%%%%%%%%%%%%%%%%%%%%

\documentclass[11pt,a4paper]{amsart}
\usepackage{xr-hyper}
\newcounter{chapter}

\usepackage{amssymb}
\usepackage{hyperref}
\usepackage{listings}
\usepackage{mathtools}
\usepackage{enumitem}

% Fuzz -------------------------------------------------------------------
\hfuzz2pt % Don't bother to report over-full boxes if over-edge is < 2pt
% Line spacing -----------------------------------------------------------

%%%%%%%%%fin pegada%%%%%%%%%%%%%%%%%%5

\usepackage[T1]{fontenc}
\usepackage[spanish]{babel}
\usepackage{graphicx}
\usepackage{tikz}
\usetikzlibrary{babel}

\setlength{\textwidth}{15.5cm} 
\setlength{\textheight}{24cm}
\setlength{\oddsidemargin}{0cm}
\setlength{\evensidemargin}{0cm}


\parskip 3pt

%%
%% OPERADORES
%%
\renewcommand{\Re}{\operatorname{Re}}
\newcommand{\tr}{\operatorname{tr}}
\newcommand{\id}{\operatorname{id}}
\def\div{\operatorname {\mathrm{div}}}
\def\Dom{\operatorname {\mathrm{Dom}}}
\def\argmin{\operatorname {\mathrm{arg\, min}}}
\def\argmax{\operatorname {\mathrm{arg\, max}}}
\def\gen{\operatorname {\mathrm{gen}}}
\def\Im{\operatorname {\mathrm{Im}}}
\def\Nu{\operatorname {\mathrm{Nu}}}
\def\sop{\operatorname {\mathrm{sop}}}
\def\diam{\operatorname {\mathrm{diam}}}
\def\dist{\operatorname {\mathrm{dist}}}
\def\var{\operatorname {\mathrm{Var}}}
\def\codim{\operatorname {\mathrm{codim}}}
\def\pint{\operatorname {--\!\!\!\!\!\int\!\!\!\!\!--}}

%%
%% EPSILON
%%
\newcommand{\ve}{\varepsilon}

%%
%% LETRAS MATHBB
%%
\newcommand{\R}{{\mathbb R}}
\newcommand{\Z}{{\mathbb Z}}
\newcommand{\C}{{\mathbb C}}
\newcommand{\Q}{{\mathbb Q}}
\newcommand{\N}{{\mathbb N}}
\newcommand{\E}{{\mathbb E}}
\newcommand{\V}{{\mathbb V}}
\renewcommand{\P}{{\mathbb P}}
\newcommand{\II}{{\mathbb I}}
\newcommand{\OO}{{\mathbb O}}
\newcommand{\Sb}{{\mathbb S}}
\newcommand{\Tb}{{\mathbb T}}

%%
%% LETRAS CALIFGRAFICAS
%%
\newcommand{\F}{{\mathcal F}}
\newcommand{\Hc}{{\mathcal H}}
\newcommand{\Tc}{{\mathcal T}}
\newcommand{\Ec}{{\mathcal E}}
\newcommand{\Vc}{{\mathcal V}}
\newcommand{\A}{{\mathcal A}}
\newcommand{\J}{{\mathcal J}}
\newcommand{\G}{{\mathcal G}}
\newcommand{\U}{{\mathcal U}}
\newcommand{\B}{{\mathcal B}}
\newcommand{\Sw}{{\mathcal S}}
\newcommand{\D}{{\mathcal D}}
\newcommand{\M}{{\mathcal M}}
\newcommand{\LL}{{\mathcal L}}

%%
%% LETRAS MATHBF
%%
\newcommand{\ab}{{\mathbf a}}
\newcommand{\bb}{{\mathbf b}}
\newcommand{\cc}{{\mathbf c}}
\newcommand{\ee}{{\mathbf e}}
\newcommand{\ff}{{\mathbf f}}
\newcommand{\gb}{{\mathbf g}}
\newcommand{\hh}{{\mathbf h}}
\newcommand{\n}{{\mathbf n}}
\newcommand{\mm}{{\mathbf m}}
\newcommand{\pp}{{\mathbf p}}
\newcommand{\qq}{{\mathbf q}}
\newcommand{\sbf}{{\mathbf s}}
\newcommand{\tb}{{\mathbf t}}
\newcommand{\uu}{{\mathbf u}}
\newcommand{\vv}{{\mathbf v}}
\newcommand{\ww}{{\mathbf w}}
\newcommand{\xx}{{\mathbf x}}
\newcommand{\yy}{{\mathbf y}}
\newcommand{\zz}{{\mathbf z}}
\newcommand{\I}{{\mathbf 1}}
\newcommand{\Pb}{{\mathbf P}}
\newcommand{\XX}{{\mathbf X}}
\newcommand{\YY}{{\mathbf Y}}

%%
%% LETRAS MATHFRAK
%%

\newcommand{\Gk}{{\mathfrak G}}


%%
%% TEOREMAS
%%
\newtheorem{teo}{Teorema}[section]
\newtheorem{lema}[teo]{Lema}
\newtheorem{prop}[teo]{Proposición}
\newtheorem{cor}[teo]{Corolario}

%%
%% REMARK
%%
\theoremstyle{remark}
\newtheorem{obs}[teo]{Observación}

%%
%% DEFINICION
%%
\theoremstyle{definition}
\newtheorem{defi}[teo]{Definición}
\newtheorem{ejem}[teo]{Ejemplo}
\newtheorem{ejer}[teo]{Ejercicio}
\newtheorem{nota}[teo]{Notación}
\newtheorem{ejercicio}{Ejercicio}[chapter]


%%
%% NUMERACION
%%
\numberwithin{subsection}{section}
\numberwithin{equation}{section}

%%
%% TITULO
%%


\externaldocument[N-]{../notas/Notas-TO}
\newcommand{\cuadernolink}[2]{\noindent\textbf{Cuaderno:} \emph{#1} (\href{https://colab.research.google.com/github/jfbonder/transporte-optimo/blob/main/cuadernos/#2}{abrir en Colab}).\par\medskip}
\pagestyle{plain}

\begin{document}
\renewcommand{\thechapter}{3}
\begin{center}
{\large\sc Teoría de Transporte Óptimo}\\[2pt]
{\small 2do cuatrimestre 2026 --- Julián Fernández Bonder}\\[10pt]
{\LARGE\bf Práctica 3}\\[4pt]
{\large Un breve repaso de dualidad en Programación Lineal}
\end{center}
\bigskip
\noindent{\small Los ejercicios son los de la sección de ejercicios del Capítulo~3 de las notas del curso, con la misma numeración. Las referencias a teoremas, proposiciones y secciones remiten a las notas (\url{https://github.com/jfbonder/transporte-optimo}).}
\medskip


\begin{ejercicio}[Duales de distintas formas]\label{ejer.PL.formas}
\begin{enumerate}[label=(\alph*)]
\item Escribir el problema dual de $\min\{c^Tx:\ Ax\ge b,\ x\ge0\}$ y el de
$\min\{c^Tx:\ Ax=b\}$ (sin restricción de signo), reduciéndolos primero a la
forma estándar $\min\{c^Tx:Ax=b,\,x\ge0\}$ mediante variables de holgura y
la descomposición $x=x^+-x^-$.
\item Probar que el dual del dual es el primal.
\item Dar un ejemplo en $\R$ (una variable) en el que el primal no sea
factible y el dual no sea acotado, y otro en el que ninguno de los dos sea
factible. ¿Contradicen estos ejemplos el teorema de dualidad fuerte?
\end{enumerate}
\end{ejercicio}

\begin{ejercicio}[Holgura complementaria]\label{ejer.PL.holgura}
Sean $x$ factible para el primal $\min\{c^Tx:Ax=b,\,x\ge0\}$ e $y$ factible
para su dual $\max\{b^Ty:A^Ty\le c\}$.
\begin{enumerate}[label=(\alph*)]
\item Probar que $c^Tx-b^Ty=x^T(c-A^Ty)$ y que ambos son óptimos si y sólo
si $x_j\,(c-A^Ty)_j=0$ para todo $j$.
\item Traducirlo al problema de transporte discreto de la
Sección~\ref{N-kant.discreto}: $\pi$ y $(\varphi,\psi)$ factibles son ambos
óptimos si y sólo si
\[
\pi_{ij}>0\ \Longrightarrow\ \varphi_i+\psi_j=c_{ij}.
\]
Ésta es la versión discreta de la \emph{condición de saturación} que
reaparecerá en los Capítulos~\ref{N-cap.dual.discreto} y~\ref{N-cap.dual}.
\end{enumerate}
\end{ejercicio}

\begin{ejercicio}[Un problema $2\times2$ a mano]\label{ejer.PL.2x2}
Sean $a=(0{,}4,\ 0{,}6)$, $b=(0{,}5,\ 0{,}5)$ y
$c=\bigl(\begin{smallmatrix}1&3\\2&1\end{smallmatrix}\bigr)$.
\begin{enumerate}[label=(\alph*)]
\item Describir $\Pi(a,b)$ como un segmento (tiene un solo grado de
libertad) y hallar el plan óptimo y el costo óptimo.
\item Resolver el problema dual y verificar la dualidad fuerte y la holgura
complementaria del Ejercicio~\ref{ejer.PL.holgura}.
\item Probar que el par óptimo $(\varphi,\psi)$ es único salvo la
transformación $(\varphi+s,\psi-s)$, $s\in\R$, e interpretar esta
invariancia.
\end{enumerate}
\end{ejercicio}

\begin{ejercicio}[Vértices del politopo de transporte]\label{ejer.PL.vertices}
Sean $a\in\R^N$, $b\in\R^M$ con entradas positivas y
$\sum a_i=\sum b_j=1$, y $\Pi(a,b)$ el politopo de planes discretos.
\begin{enumerate}[label=(\alph*)]
\item Probar que $\Pi(a,b)$ es un politopo compacto y convexo de dimensión
$(N-1)(M-1)$.
\item Probar que todo punto extremal $\pi$ de $\Pi(a,b)$ tiene a lo sumo
$N+M-1$ entradas positivas. \emph{Sugerencia:} el grafo bipartito con
aristas $\{(i,j):\pi_{ij}>0\}$ no puede contener ciclos (si los contiene,
perturbar como en el Ejercicio~\ref{N-ejer.biestocasticas}), y un bosque con
$N+M$ vértices tiene a lo sumo $N+M-1$ aristas.
\item Deducir que el problema de transporte discreto siempre admite un plan
óptimo con a lo sumo $N+M-1$ entradas positivas. (Esta cota se usa en el
Capítulo~\ref{N-cap.Sinkhorn} para contrastar con el plan entrópico, que
tiene todas sus entradas positivas.)
\end{enumerate}
\end{ejercicio}

\begin{ejercicio}[El problema de asignación]\label{ejer.PL.asignacion}
Sean $N=M$ y $a=b=\frac1N(1,\dots,1)$.
\begin{enumerate}[label=(\alph*)]
\item Probar, usando que un programa lineal alcanza su mínimo en un vértice
y el teorema de Birkhoff (Ejercicio~\ref{N-ejer.biestocasticas}), que
\[
\min_{\pi\in\Pi(a,b)}\sum_{ij}c_{ij}\pi_{ij}
=\frac1N\min_{\sigma\in S_N}\sum_ic_{i\sigma(i)}.
\]
Es decir: para pesos uniformes, el problema de Kantorovich discreto coincide
con el problema de Monge (una asignación óptima), aunque el primero tiene
$N^2$ variables continuas y el segundo $N!$ candidatos.
\item Escribir el dual del problema de asignación y probar que una
permutación $\sigma$ es óptima si y sólo si existen $\varphi,\psi\in\R^N$
con $\varphi_i+\psi_j\le c_{ij}$ para todo $i,j$ y
$\varphi_i+\psi_{\sigma(i)}=c_{i\sigma(i)}$ para todo $i$. Interpretar
$\varphi_i$ y $\psi_j$ como precios (problema de asignación de
Koopmans--Beckmann).
\item Para $c_{ij}=|x_i-y_j|^2$ con $x_1<\dots<x_N$ e $y_1<\dots<y_N$ en
$\R$, probar directamente que $\sigma=\id$ es óptima (intercambiar dos
asignaciones cruzadas no aumenta el costo). Éste es el germen del
Capítulo~\ref{N-cap.1D}.
\end{enumerate}
\end{ejercicio}

\end{document}
